\documentclass{report}
\usepackage[utf8]{inputenc}
\usepackage[T1]{fontenc}
\usepackage[french]{babel}
\usepackage{geometry}
\usepackage{fancyhdr}
\setlength{\headheight}{13pt}
\usepackage{lastpage}
\usepackage{titlesec}
\usepackage{listings}
\usepackage{enumitem}
\usepackage{color}
\usepackage[table]{xcolor}
\usepackage{soul}
\usepackage{graphicx}
\usepackage{amsmath}
\usepackage{amssymb}
\usepackage{hyperref}
\hypersetup{hidelinks}
\usepackage{tikz}
\usetikzlibrary{positioning, shapes.geometric, arrows.meta, calc, fit, backgrounds, shadows}
\usepackage{rotating}
\usepackage{longtable}
\usepackage{array}

\newcommand{\DocName}{OPTIMISATION DE MODELES NEURONAUX A OSCILLATEURS COUPLÉS - LIVRABLE 3 : ÉTAT DE L’ART SOUS FORME DE BENCHMARK}
\newcommand{\DocAuthors}{\bsc{Da Silva} Tristan - \bsc{Battaglia} Charles-Antoine \\ \bsc{Botté-Magalhaes} Jules - \bsc{Gauthier} Keanu - \bsc{Fraysse} Ianis}

\definecolor{dkgreen}{rgb}{0,0.6,0}
\definecolor{gray}{rgb}{0.5,0.5,0.5}
\definecolor{mauve}{rgb}{0.58,0,0.82}
\definecolor{addrred}{rgb}{0.93,0.43,0.43}
\definecolor{commentgray}{rgb}{0.45,0.45,0.45}

% Palette du benchmark
\definecolor{bmBest}{HTML}{B5EAD7}
\definecolor{bmGood}{HTML}{C7E9C0}
\definecolor{bmMid}{HTML}{FFE699}
\definecolor{bmLow}{HTML}{FFC0CB}

\lstset{frame=tb,
    aboveskip=3mm,
    belowskip=3mm,
    showstringspaces=false,
    columns=flexible,
    basicstyle={\small\ttfamily},
    numbers=none,
    numberstyle=\tiny\color{gray},
    keywordstyle=\color{blue},
    commentstyle=\color{commentgray},
    breaklines=true,
    breakatwhitespace=true,
    tabsize=3,
   }

\lstset{breaklines=true}

\geometry{a4paper, margin=1in}

% Suppression de l'affichage du mot "Chapitre"
\titleformat{\chapter}[display]
    {\normalfont\bfseries\LARGE}
    {}
    {0pt}
    {}

\definecolor{Light}{gray}{.90}
\sethlcolor{Light}

\let\OldTexttt\texttt
\renewcommand{\texttt}[1]{\OldTexttt{\hl{#1}}}

\renewcommand{\thechapter}{\arabic{chapter}}
\renewcommand{\thesection}{\thechapter.\arabic{section}}
\renewcommand{\thesubsection}{\thesection.\arabic{subsection}}

\fancyhf{}
\fancyhead[L]{\DocName}
\fancyfoot[L]{\DocAuthors}
\fancyfoot[R]{\thepage\ / \pageref{LastPage}}
\renewcommand{\headrulewidth}{0 pt}
\renewcommand{\footrulewidth}{0.2pt}
\pagestyle{fancy}
\fancypagestyle{plain}{
    \fancyhf{}
    \fancyfoot[R]{\thepage\ / \pageref{LastPage}}
    \fancyfoot[L]{\DocAuthors}
    \fancyhead[L]{\DocName}
}

\graphicspath{ {./img} }

\title{\DocName}
\author{\DocAuthors}
\date{\textbf{JUNIA - ISEN - AP4} \url{https://github.com/CharlesBta/Modele-Neuronaux-RD.git}}

\begin{document}
\maketitle

\chapter*{Lexique}
\addcontentsline{toc}{chapter}{Lexique}

\begin{tabular}{@{}p{3.9cm}p{11cm}@{}}
    \textbf{Backend} & Moteur d'exécution d'un simulateur ; Brian2CUDA ajoute par exemple un backend \textbf{GPU} au simulateur Brian. \\[6pt]
    \textbf{Génération de code} & Production automatique du code \textbf{GPU} à partir d'une description de modèle écrite en langage de haut niveau (\textbf{Python} ou \textbf{Julia}). \\[6pt]
    \textbf{JAX} & Bibliothèque \textbf{Python} de calcul différentiable avec vectorisation automatique (vmap). \\[6pt]
    \textbf{Jonction communicante} & \emph{Gap junction} : couplage électrique direct et bidirectionnel entre deux cellules voisines. \\[6pt]
    \textbf{Julia} & Langage de programmation scientifique haut niveau, proche du \textbf{C++} en performance. \\[6pt]
    \textbf{MPI} & \emph{Message Passing Interface} : norme d'échange de messages entre machines, utilisée pour répartir une simulation sur plusieurs nœuds de calcul. \\[6pt]
    \textbf{NEST} & Simulateur de référence pour les grandes populations de neurones simplifiés. \\[6pt]
    \textbf{NVLink} & Liaison matérielle à très haut débit entre \textbf{GPU} \textbf{NVIDIA}. \\[6pt]
    \textbf{PyTorch} & Bibliothèque \textbf{Python} d'apprentissage automatique, aussi utilisée pour du calcul \textbf{GPU} généraliste. \\[6pt]
    \textbf{vmap} & Opération de vectorisation automatique de \textbf{JAX} et \textbf{PyTorch} : une même fonction s'applique à un grand nombre d'éléments en parallèle. \\
\end{tabular}

\tableofcontents

\chapter{Introduction}

Ce document est le livrable L3 du projet : l'état de l'art sous forme de benchmark. La note de clarification (L2) a posé le besoin : savoir si les serveurs \textbf{GPU} \textbf{RTX~5090} de Junia permettent de simuler des réseaux de neurones de \textbf{Hodgkin-Huxley} couplés nettement plus grands que sur \textbf{CPU}, dans un temps acceptable et sans dérive numérique. Ce document cherche ce que la littérature dit déjà de cette question, pour préparer nos propres mesures.

Nous avons retenu cinq travaux parus entre 2008 et 2025, cherchés dans des sources reconnues : \textbf{IEEE Xplore}, PubMed Central, ScienceDirect, l'\textbf{ACM} Digital Library et arXiv. Les critères de sélection : simuler des réseaux de neurones (de préférence à conductance, comme Hodgkin-Huxley) ou des systèmes d'\textbf{EDO} sur \textbf{GPU}, avec des mesures de performance publiées. Chacun fait l'objet d'une fiche d'analyse à six questions (chapitre 1) ; leur synthèse donne le benchmark comparatif du chapitre 2.

D'autres sources complètent cette analyse sans faire l'objet d'une fiche : la fiche technique de la \textbf{RTX~5090}, notre matériel cible \cite{techpowerup} ; une simulation de réseaux fortement connectés sur \textbf{GPU}, notre pire cas \cite{golosio2021} ; la parallélisation sur \textbf{GPU} du solveur des équations d'un neurone multi-compartiments \cite{vooturi2018} ; et la multiplication de matrices creuses sur \textbf{GPU} \cite{bell2009}, la forme sous laquelle notre connectome sera stocké. Toutes sont citées en bibliographie.

\chapter{Fiches d'analyse des articles}

\section{Scalable Parallel Programming with CUDA (2008)}

Cet article fondateur, écrit par des ingénieurs \textbf{NVIDIA} et un universitaire \cite{nickolls2008}, explique le modèle de programmation \textbf{CUDA} au moment où les \textbf{CPU} deviennent multicœurs et les \textbf{GPU} massivement parallèles.

\medskip
\begin{tabular}{|p{3.2cm}|p{11.2cm}|}
\hline
\textbf{Pourquoi ?} & Les processeurs grand public gagnent désormais en nombre de cœurs plutôt qu'en fréquence : le matériel devient parallèle, et les logiciels classiques ne le sont pas. Le problème posé : écrire des programmes qui passent à l'échelle quel que soit le nombre de cœurs. \\ \hline
\textbf{Comment ?} & Le modèle \textbf{CUDA} repose sur trois abstractions : une hiérarchie de groupes de fils d'exécution, des mémoires partagées et des barrières de synchronisation. Le programmeur répartit le travail en blocs de fils sans connaître le nombre de cœurs : le matériel s'occupe de l'ordonnancement. \\ \hline
\textbf{Avec quoi ?} & Cartes \textbf{GPU} \textbf{NVIDIA} et kit de développement \textbf{CUDA}, illustrés par des exemples d'applications scientifiques parallèles. \\ \hline
\textbf{Quels résultats ?} & Un même programme \textbf{CUDA} s'exécute sur des \textbf{GPU} de quelques dizaines ou de plusieurs centaines de cœurs sans être réécrit : le parallélisme suit le matériel. \\ \hline
\textbf{Quelles limites ?} & Article de présentation de 2008 : matériel et outils datés, pas de mesure sur des \textbf{EDO} biologiques, et le programmeur travaille toujours à un niveau bas (gestion fine de la mémoire et des fils). \\ \hline
\textbf{Pour notre sujet ?} & C'est le socle conceptuel du projet : paralléliser Hodgkin-Huxley voudra dire répartir les neurones sur des fils d'exécution et organiser la mémoire partagée pour le couplage synaptique. \\ \hline
\end{tabular}

\newpage

\section{GeNN : a code generation framework for accelerated brain simulations (2016)}

GeNN (\emph{GPU-enhanced Neuronal Networks}) est une bibliothèque open source développée à l'université de Sussex pour mettre les \textbf{GPU} à portée des neuroscientifiques \cite{yavuz2016}.

\medskip
\begin{tabular}{|p{3.2cm}|p{11.2cm}|}
\hline
\textbf{Pourquoi ?} & Simuler des réseaux de neurones réalistes est lent, et les accélérateurs \textbf{GPU} restent inaccessibles à ceux qui ne savent pas programmer à bas niveau. \\ \hline
\textbf{Comment ?} & L'utilisateur décrit ses neurones et ses synapses dans une interface de haut niveau ; GeNN génère automatiquement le code \textbf{CUDA} correspondant, sans exiger de connaissances \textbf{GPU} approfondies. \\ \hline
\textbf{Avec quoi ?} & \textbf{GPU} \textbf{NVIDIA} ; bancs d'essai sur plusieurs modèles de réseaux, dont un réseau d'un million de neurones de \textbf{Hodgkin-Huxley} à conductance. \\ \hline
\textbf{Quels résultats ?} & Jusqu'à 200 fois plus rapide qu'un cœur de \textbf{CPU} pour le réseau d'un million de neurones \textbf{Hodgkin-Huxley} ; mais le gain change selon le modèle simulé. \\ \hline
\textbf{Quelles limites ?} & Gain très dépendant du modèle ; comparaison faite contre un seul cœur de \textbf{CPU}, donc une référence faible ; \textbf{NVIDIA} uniquement. \\ \hline
\textbf{Pour notre sujet ?} & Preuve directe qu'un réseau de \textbf{Hodgkin-Huxley} à conductance passe à l'échelle sur \textbf{GPU}, et première forme de l'idée de génération de code que notre projet pourra réutiliser ou comparer à des noyaux écrits à la main. \\ \hline
\end{tabular}

\section{Brian2CUDA : Flexible and Efficient Simulation of Spiking Neural Network Models on GPUs (2022)}

Brian2CUDA ajoute un backend \textbf{GPU} au simulateur Brian2, très utilisé en neurosciences computationnelles \cite{alevi2022}.

\medskip
\begin{tabular}{|p{3.2cm}|p{11.2cm}|}
\hline
\textbf{Pourquoi ?} & Les chercheurs décrivent leurs modèles en \textbf{Python} avec le simulateur Brian2, mais celui-ci n'avait pas de sortie \textbf{GPU} complète : écrire le code bas niveau soi-même demandait un savoir-faire rare. \\ \hline
\textbf{Comment ?} & Le simulateur génère lui-même le code \textbf{CUDA} à partir de la définition du modèle en \textbf{Python} : un volet pour l'intégration numérique des états des neurones, un autre pour la propagation des événements synaptiques. \\ \hline
\textbf{Avec quoi ?} & Code \textbf{CUDA} généré ; comparaisons avec le backend \textbf{CPU} de Brian sur plusieurs types de modèles, y compris neurones et synapses arbitraires et plasticité. \\ \hline
\textbf{Quels résultats ?} & Jusqu'à trois ordres de grandeur d'accélération face au backend \textbf{CPU} de Brian ; seule solution \textbf{GPU} qui couvre toutes les fonctionnalités de Brian2. \\ \hline
\textbf{Quelles limites ?} & Lié au cadre Brian2 : on doit écrire ses modèles dans ce simulateur ; gain variable selon les modèles ; un seul \textbf{GPU}. \\ \hline
\textbf{Pour notre sujet ?} & Démontre qu'on peut accélérer sans écrire de \textbf{CUDA} à la main : le modèle décrit en \textbf{Python} est traduit. Notre socle \textbf{Python}/SciPy pourrait suivre la même voie : compiler les équations plutôt que réécrire le solveur. \\ \hline
\end{tabular}
\newpage
\section{Automated Translation and Accelerated Solving of Differential Equations on Multiple GPU Platforms (2024)}

Cet article, publié dans Computer Methods in Applied Mechanics and Engineering \cite{utkarsh2024}, cherche à résoudre beaucoup d'\textbf{EDO} sur \textbf{GPU} sans écrire de code compliqué, et quel que soit le constructeur de la carte.

\medskip
\begin{tabular}{|p{3.2cm}|p{11.2cm}|}
\hline
\textbf{Pourquoi ?} & Résoudre beaucoup d'\textbf{EDO} sur \textbf{GPU} oblige normalement à écrire du code \textbf{CUDA} à la main, et ce code ne fonctionne que sur les cartes \textbf{NVIDIA}. \\ \hline
\textbf{Comment ?} & Le solveur d'\textbf{EDO} du langage \textbf{Julia} traduit tout seul le code pour la carte : on écrit ses équations comme d'habitude, et elles s'exécutent sur le \textbf{GPU} sans rien changer au modèle. \\ \hline
\textbf{Avec quoi ?} & Des \textbf{GPU} de quatre constructeurs (\textbf{NVIDIA}, AMD, Intel, Apple), et des comparaisons avec du \textbf{CUDA-C++} écrit à la main puis avec \textbf{JAX} et \textbf{PyTorch}. \\ \hline
\textbf{Quels résultats ?} & Aussi rapide que du \textbf{CUDA-C++} optimisé à la main, et 20 à 100 fois plus rapide que la solution de \textbf{JAX} et \textbf{PyTorch}. Le même code tourne sur toutes les cartes testées. \\ \hline
\textbf{Quelles limites ?} & Les \textbf{EDO} de l'article sont indépendantes les unes des autres : pas de synapses ni de connectome, donc pas de couplage irrégulier comme dans notre cas. Et il faut travailler en \textbf{Julia}. \\ \hline
\textbf{Pour notre sujet ?} & Notre réseau est justement un gros système d'\textbf{EDO}. L'article montre qu'on peut éviter d'écrire le \textbf{CUDA} nous-mêmes et ne pas dépendre d'un seul constructeur. Reste à vérifier si ça marche aussi bien quand les équations sont couplées entre elles. \\ \hline
\end{tabular}

\section{SimHH : A Versatile, Multi-GPU Simulator for Extended Hodgkin-Huxley Networks (2025)}

SimHH est le travail le plus récent et le plus proche de notre cas : un simulateur de réseaux de \textbf{Hodgkin-Huxley} étendus, distribué sur des clusters de calcul \cite{engelen2025}.

\medskip
\begin{tabular}{|p{3.2cm}|p{11.2cm}|}
\hline
\textbf{Pourquoi ?} & Les réseaux de \textbf{Hodgkin-Huxley} étendus (multi-compartiments, jonctions communicantes) dépassent ce qu'une seule machine peut simuler en un temps raisonnable. \\ \hline
\textbf{Comment ?} & Répartition de la simulation sur plusieurs nœuds de calcul avec \textbf{MPI} (OpenMPI), avec deux configurations : partager tous les potentiels entre nœuds, ou ne partager que les potentiels utiles ; \textbf{CUDA} et \textbf{NVLink} pour le calcul, les noyaux tournant en parallèle des communications. \\ \hline
\textbf{Avec quoi ?} & Clusters multi-\textbf{GPU} allant jusqu'à 256 nœuds ; modèles étendus dont le noyau inféro-olivaire ; comparaisons avec des \textbf{CPU}, un \textbf{FPGA} et le simulateur \textbf{NEST}. \\ \hline
\textbf{Quels résultats ?} & Réseaux de plus de 10 millions de cellules ; environ 150$\times$ plus rapide qu'un \textbf{CPU} mono-fil, 10$\times$ qu'un \textbf{CPU} à 128 fils, 10$\times$ qu'un \textbf{FPGA}, et 7$\times$ que \textbf{NEST} sur 32 nœuds pour un réseau de 94\,720 neurones à jonctions communicantes. \\ \hline
\textbf{Quelles limites ?} & Exige un cluster de calcul (plusieurs machines reliées) et \textbf{MPI} : une infrastructure bien au-delà d'un serveur universitaire ; la configuration de répartition est délicate à régler. \\ \hline
\textbf{Pour notre sujet ?} & État de l'art le plus proche : \textbf{Hodgkin-Huxley}, comparaisons honnêtes (\textbf{CPU} multi-fils, \textbf{FPGA}, simulateurs) et preuve qu'un seul \textbf{GPU} finit par saturer. Nos mesures se placeront juste avant cette étape multi-\textbf{GPU}. \\ \hline
\end{tabular}

\chapter{Benchmark comparatif}

Le tableau ci-dessous compare les cinq approches sur les critères de benchmarking. Les gains de performance sont indiqués par rapport à la référence mesurée dans chaque article.

\medskip

{\setlength{\tabcolsep}{3pt}\footnotesize
\begin{tabular}{|p{1.8cm}|p{1.5cm}|p{1.5cm}|p{1.85cm}|p{1.05cm}|p{1.35cm}|p{2.1cm}|p{0.95cm}|}
\hline
\rowcolor{gray!15}
\scriptsize\textbf{Outil} & \scriptsize\textbf{Type} & \scriptsize\textbf{Technicité} & \scriptsize\textbf{Implantation} & \scriptsize\textbf{Coût} & \scriptsize\textbf{Facilité} & \scriptsize\textbf{Gains de performance} & \scriptsize\textbf{Moy.} \\ \hline

CUDA \cite{nickolls2008} & Modèle de programmation \textbf{GPU} &
\cellcolor{bmBest}Très avancé &
\cellcolor{bmBest}Standard en production &
\cellcolor{bmBest}Gratuit &
\cellcolor{bmLow}Complexe &
Sans mesure &
\cellcolor{bmGood}+ \\ \hline

GeNN \cite{yavuz2016} & Génération de code \textbf{CUDA} &
\cellcolor{bmGood}Avancé &
\cellcolor{bmGood}Utilisé en recherche &
\cellcolor{bmBest}Gratuit &
\cellcolor{bmGood}Assez facile &
\cellcolor{bmGood}$\times$200 vs 1 cœur \textbf{CPU} &
\cellcolor{bmGood}+ \\ \hline

Brian2CUDA \cite{alevi2022} & Backend \textbf{GPU} de Brian2 &
\cellcolor{bmMid}Modéré &
\cellcolor{bmMid}Recherche, récent &
\cellcolor{bmBest}Gratuit &
\cellcolor{bmBest}Très facile &
\cellcolor{bmBest}$\times$1000 vs Brian \textbf{CPU} &
\cellcolor{bmGood}+ \\ \hline

Solveur \textbf{Julia} \cite{utkarsh2024} & Solveur d'\textbf{EDO} multi-\textbf{GPU} &
\cellcolor{bmGood}Avancé &
\cellcolor{bmGood}Écosystème actif &
\cellcolor{bmBest}Gratuit &
\cellcolor{bmGood}Facile &
\cellcolor{bmBest}$\times$20--100 vs \textbf{JAX}/\allowbreak\textbf{PyTorch} &
\cellcolor{bmGood}+ \\ \hline

SimHH \cite{engelen2025} & Multi-\textbf{GPU} (\textbf{MPI}) &
\cellcolor{bmBest}Très avancé &
\cellcolor{bmMid}Prototype d'une équipe &
\cellcolor{bmLow}Exige un cluster &
\cellcolor{bmLow}Complexe &
\cellcolor{bmBest}$\times$150 vs \textbf{CPU} mono-fil &
\cellcolor{bmMid}- \\ \hline

\end{tabular}}

\medskip

\begin{tabular}{|p{1.5cm}|p{6cm}|}
\hline
\textbf{Couleur} & \textbf{Signification} \\ \hline
\cellcolor{bmBest} & \textbf{++} Meilleur cas \\ \hline
\cellcolor{bmGood} & \textbf{+} Bon cas \\ \hline
\cellcolor{bmMid} & \textbf{-} Cas plutôt moyen \\ \hline
\cellcolor{bmLow} & \textbf{--} Cas moyen \\ \hline
& Non concerné \\ \hline
\end{tabular}

\begin{thebibliography}{9}
\addcontentsline{toc}{chapter}{Bibliographie}

\bibitem{nickolls2008} J.~Nickolls, I.~Buck, M.~Garland, K.~Skadron, \og Scalable Parallel Programming with CUDA \fg, \emph{ACM Queue}, vol.~6, n\textsuperscript{o}~2, p.~24--32, mars 2008. \href{https://doi.org/10.1145/1365490.1365500}{DOI : 10.1145/1365490.1365500}.

\bibitem{yavuz2016} E.~Yavuz, J.~Turner, T.~Nowotny, \og GeNN : a code generation framework for accelerated brain simulations \fg, \emph{Scientific Reports}, vol.~6, article 18854, 2016. \href{https://doi.org/10.1038/srep18854}{DOI : 10.1038/srep18854}.

\bibitem{alevi2022} D.~Alevi, M.~Stimberg, H.~Sprekeler, K.~Obermayer, M.~Augustin, \og Brian2CUDA : Flexible and Efficient Simulation of Spiking Neural Network Models on GPUs \fg, \emph{Frontiers in Neuroinformatics}, vol.~16, article 883700, 2022. \href{https://doi.org/10.3389/fninf.2022.883700}{DOI : 10.3389/fninf.2022.883700}.

\bibitem{utkarsh2024} U.~Utkarsh, V.~Churavy, Y.~Ma, T.~Besard, P.~Srisuma, T.~Gymnich, A.~R.~Gerlach, A.~Edelman, G.~Barbastathis, R.~D.~Braatz, C.~Rackauckas, \og Automated Translation and Accelerated Solving of Differential Equations on Multiple GPU Platforms \fg, \emph{Computer Methods in Applied Mechanics and Engineering}, vol.~419, 2024. \href{https://doi.org/10.1016/j.cma.2023.116591}{DOI : 10.1016/j.cma.2023.116591}. Prépublication : \href{https://arxiv.org/abs/2304.06835}{arXiv:2304.06835}.

\bibitem{engelen2025} M.~C.~W.~Engelen, J.-H.~L.~F.~Betting, C.~Strydis, \og SimHH : A Versatile, Multi-GPU Simulator for Extended Hodgkin-Huxley Networks \fg, \emph{IEEE Access}, vol.~13, 2025. \href{https://doi.org/10.1109/ACCESS.2025.3550444}{DOI : 10.1109/ACCESS.2025.3550444}.

\bibitem{golosio2021} B.~Golosio, G.~Tiddia, C.~De~Luca, E.~Pastorelli, F.~Simula, P.~S.~Paolucci, \og Fast Simulations of Highly-Connected Spiking Cortical Models Using GPUs \fg, \emph{Frontiers in Computational Neuroscience}, vol.~15, article 627620, 2021. \href{https://doi.org/10.3389/fncom.2021.627620}{DOI : 10.3389/fncom.2021.627620}.

\bibitem{vooturi2018} D.~T.~Vooturi, K.~Kothapalli, U.~S.~Bhalla, \og Parallelizing Hines Matrix Solver in Neuron Simulations on GPU \fg, \emph{2017 IEEE 24th International Conference on High Performance Computing (HiPC)}, p.~388--397, 2018. \href{https://doi.org/10.1109/HiPC.2017.00051}{DOI : 10.1109/HiPC.2017.00051}.

\bibitem{bell2009} N.~Bell, M.~Garland, \og Implementing sparse matrix-vector multiplication on throughput-oriented processors \fg, \emph{Proceedings of the Conference on High Performance Computing, Networking, Storage and Analysis (SC'09)}, ACM, 2009. \href{https://doi.org/10.1145/1654059.1654078}{DOI : 10.1145/1654059.1654078}.

\bibitem{techpowerup} \og NVIDIA GeForce RTX 5090 \fg, fiche technique, TechPowerUp. \href{https://www.techpowerup.com/gpu-specs/geforce-rtx-5090.c4216}{www.techpowerup.com/gpu-specs/geforce-rtx-5090.c4216}.

\end{thebibliography}

\end{document}
